\section{模型假设}

为保证模型可计算性、可验证性，并避免作出超出题目数据支持范围的生理推断，本文作如下基本假设：

\begin{enumerate}[label=\arabic*.]
    \item \textbf{数据有效性假设。} 题目提供的采样顺序、时间戳、视觉提示标记以及F3、Fz、F4原始脑电记录均真实有效，采集系统在同一记录内保持稳定，不考虑人为篡改和大规模硬件故障。
    \item \textbf{相对幅值可比假设。} 在同一实验数据来源内，各试次由同一记录系统获得，因而可在原始记录单位下比较相对幅值、波形形态和时间位置；由于缺少可靠标定信息，本文不擅自将幅值换算为$\mu\mathrm{V}$。
    \item \textbf{刺激前基线假设。} 短时间刺激前区间可近似代表该试次的局部基线，基线校正主要消除直流偏置和缓慢漂移，不改变刺激后的相对动力学结构。
    \item \textbf{事件锁定假设。} 与视觉提示直接相关的神经响应在重复试次中存在一定的事件锁定共同结构，而部分眼动、肌电和随机背景活动与提示事件不同步，因此跨试次统计能够提高视觉诱发成分的可见度。
    \item \textbf{有限空间可辨识假设。} F3、Fz、F4仅提供三个头皮观测，不足以唯一反演LGN、海马或完整皮层的真实三维电流源。因此问题二、三只建立低维等效神经状态与头皮观测之间的映射，不将模型隐状态解释为唯一的解剖定位结果。
    \item \textbf{任务结构共享假设。} 左、右视觉刺激共享大部分基础视觉通路和观测机制，其差异主要通过刺激空间配置、输入编码或部分神经群体状态体现，从而允许在共享动力学框架下比较左右条件。
    \item \textbf{行为与认知非等价假设。} 目标应答可受动作执行、误操作或反应延迟影响，因此行为正确与否不直接等同于认知识别成功与否；问题三中行为通道仅作为外部辅助观测，不作为潜在认知状态的唯一真值标签。
\end{enumerate}
